% =============================================================================
% Table — Conceptual Distinctions in the Identity–Determination Relation
% Book I — Ontology of Viability
% Chapter 01 — Identity
% Section 03 — Identity and Determination
%
% Asset ID:
% VS.B01.P01.C01.S03.SS10.P001.TL01
%
% Storage:
% 4_ASSETS/tables/
% =============================================================================

\begin{table}[htbp]
\centering
\caption{Conceptual Distinctions in the Identity--Determination Relation}
\label{tab:identity-determination-conceptual-distinctions}

\small
\renewcommand{\arraystretch}{1.3}

\begin{tabular}{p{0.22\textwidth} p{0.34\textwidth} p{0.34\textwidth}}
\hline
\textbf{Conceptual Dimension}
&
\textbf{Meaning}
&
\textbf{Relation to Identity}
\\
\hline

General Determination
&
The ontological condition through which an existent is determinate rather than wholly unspecified.
&
Establishes the general context in which identity must be considered, but does not identify the determinations constitutive of identity.
\\
\hline

Internal Determination
&
A determination considered in relation to the constitution of the existent itself.
&
May participate in constituting identity, but internality alone does not establish constitutive significance.
\\
\hline

Relational Determination
&
A determination arising through or constituted within relations involving the existent and other entities or conditions.
&
May participate in identity constitution; relational character does not establish either inclusion or exclusion.
\\
\hline

Structural Determination
&
The organization, configuration, and interdependence of determinations within an existent.
&
May contribute to identity constitution, but not every structural feature is necessarily constitutive.
\\
\hline

Material Determination
&
A determination associated with the material constitution of an existent.
&
May be indispensable to the identity of some existents, but material constitution is not a universal definition of identity.
\\
\hline

Identity-Constitutive Determination
&
A determination whose ontological role participates in constituting an existent as the existent that it is.
&
Directly concerns the constitutive relation under investigation; its attribution requires ontological justification.
\\
\hline

Descriptive Determination
&
A determination as articulated, selected, or represented within a description of an existent.
&
Its descriptive identification does not by itself establish its ontological status or constitutive significance.
\\
\hline

\end{tabular}

\vspace{0.5em}

\begin{minipage}{0.92\textwidth}
\footnotesize
\textit{Note:} These dimensions are analytical distinctions rather
than an exhaustive or mutually exclusive taxonomy of ontological
determinations. A particular determination may belong to more
than one analytical dimension. Identity-constitutive significance
is a separate question and cannot be inferred solely from
classification within any one dimension.
\end{minipage}

\end{table}